\scanpage{15}
\mspage{7}
Considérons le groupe $\mathfrak{D}$ en tant que groupe ``abstrait'', il est canoniquement extension de $\mathfrak{D}_{\mathrm{ab}}=W$ ($\simeq\mathbb{F}_2\times\mathbb{F}_2$ non canoniquement) par $\mathbb{F}_2$. Dans $\mathfrak{D}_{\mathrm{ab}}$ il y a un élément non nul distingué, \inserted{$\varpi$}, \deleted{\unsure{image inverse}}\inserted{dont l'image} inverse est l'ens $\underline{\omega}_Q$ des deux éléments d'ordre $4$ de $\mathfrak{D}$. On voit aisément que les deux autres éléments non nuls de $\mathfrak{D}_{\mathrm{ab}}$ sont interchangés par le groupe des autom. de $\mathfrak{D}$, on trouve donc
\diagram{\begin{tikzcd}[column sep=large]
1 \arrow[r] &
\underset{\substack{\text{autom. de l'ext. $\mathfrak{D}$ de}\\
\text{$W$ par $\mathbb{F}_2$ qui}\\
\text{induisent l'identité}\\
\text{dans $W$ et $\mathbb{F}_2$}}}{\check W}
\arrow[r] & \operatorname{Aut}\mathfrak{D} \arrow[r] &
\underbrace{\pm1}_{\substack{\text{interchange}\\
\text{les les deux}\\
\text{éléments de}\\
W\setminus\{0,\varpi\}}} \arrow[r,dash] & 1
\end{tikzcd}}
\margin{\uline{NB} on a en fait un iso. can. $\check W\simeq W$ ($\forall$ vectoriel de rg $2$ sur $\mathbb{F}_2$).}
$\operatorname{Aut}\mathfrak{D}$ est pr. semi-direct de $\pm1$ par $\check W$, $\pm1$ opérant par échange des $2$ éléments non \deleted{\unsure{distin.}} privilégiés de $\check W\simeq W$. Quand on tient compte que $\mathfrak{D}$ provient d'un carré, donc il y a choix d'un élément $\underline{c}$ de $W\setminus\{0,\varpi\}$ (correspondant \deleted{\unsure{aux}} à $C=\!V_S\setminus\{1,a\}$), l'autre élément étant $\underline{d}$ (correspondant à $D=V_A\setminus\{1,a\}$).
Les autom. de $\mathfrak{D}$ qui respectent cette structure forment le sous-groupe $\check W\simeq W$ de $\operatorname{Aut}\mathfrak{D}$, formé des autom. \inserted{d'ext.} qui induisent l'identité sur le quotient $W$. Ce sont aussi les autom. intérieurs, qui forment bien le groupe $W\simeq\mathfrak{D}/\underbrace{\{1,a\}}_{\text{centre}}$. (\unsure{Comme} l'hom $\mathfrak{D}\xrightarrow{\mathrm{int}}\operatorname{Aut}\mathfrak{D}$ n'est pas injectif), on voit que la donnée du groupe $\mathfrak{D}=\operatorname{Aut}_Q$ (qui revient à la donnée de $C$ et de $D$, cf plus bas)
\deleted{\unsure{revient à} $\mathfrak{D}=\mathbb{F}_2^C\mathbin{*_{\mathbb{F}_2}}\mathbb{F}_2^D$ \unsure{)}}
% archive foliation: 14
